\documentclass{article}

\usepackage[utf8]{inputenc}     % pdfLaTeX
\usepackage[T1]{fontenc}
\usepackage[magyar]{babel}
\usepackage{amsmath}

% Set page size and margins
% Replace `letterpaper' with `a4paper' for UK/EU standard size
\usepackage[a4paper,top=2cm,bottom=2cm,left=3cm,right=3cm,marginparwidth=1.75cm]{geometry}
%for arrows
\usepackage{textcomp}

\title{Magánhangzók csoportosítása}
\author{Lengyel Zoltán Péter\thanks{Csukás Ibolya tanítása alapján}}
\date{Legutóbb frissítve: 2026. Március 27.}

\begin{document}

\maketitle
\tableofcontents
\newpage

 \section{Magánhangzók csoportosítása}
 \subsection{Időtartam szerint}
 \begin{itemize}
     \item Rövid: a, e, i, o, ö, u, ü
     \item Hosszú: á, é, í, ó, ő, ú, ű
 \end{itemize}

 \subsection{A nyelv függőleges mozgása szerint}
 \begin{itemize}
     \item Felső nyelvállású: i, í, u, ú, ü, ű
     \item Középső nyelvállású: é, o, ó, ö, ő
     \item Alsó nyelvállású: a, á, e
 \end{itemize}

\subsection{A nyelv vízszintes mozgása szerint}
\begin{itemize}
    \item Magas (elöl képzett): e, é, i, í, ö, ő
    \item Mély (hátul képzett): a á, o, ó, u, ú
\end{itemize}

\subsection{Ajakműködés szerint}
\begin{itemize}
    \item Ajakkerekítéses: a, o, ó, ö, ő, u, ú, ü, ű
    \item Ajakréses: á, e, é, i, í
\end{itemize}

\end{document}